\documentclass[10pt,a4paper]{amsart}
\usepackage[margin=19mm]{geometry}
\usepackage{amsmath,amssymb,mathtools}
\usepackage[scr=rsfs,cal=euler]{mathalfa}
\usepackage{xfrac}
\usepackage[unicode,hidelinks,bookmarks=false]{hyperref}
\newcommand{\ie}{i.e.\ }
\newcommand{\cf}{cf.\ }
\newcommand{\resp}{respectively\ }
\newcommand{\Subsection}{\subsection}
\providecommand{\nobreakdash}{\nobreak\hbox{-}}
\setlength{\emergencystretch}{2em}
\multlinegap0pt
\usepackage{amssymb}
\usepackage{mathtools}
\usepackage{bm}
\usepackage{xcolor}
\newtheorem{theorem}{Theorem}
\newtheorem{proposition}{Proposition}
\newcommand{\KL}{D_{\mathrm{KL}}}
\newcommand{\NK}{N_K}
\newcommand{\NV}{N_V}
\newcommand{\PV}{P}
\newcommand{\QVm}{Q^{m}}
\newcommand{\tr}{\operatorname{tr}}
\title{Information-Geometric Decomposition of\\ Generalization Error in Unsupervised Learning}
\author{Gilhan Kim}
\date{Source: arXiv:2604.12340v1 (CC BY 4.0)}
\begin{document}
\maketitle

\noindent\emph{Source and licence.} Information-Geometric Decomposition of\\ Generalization Error in Unsupervised Learning. Version arXiv:2604.12340v1 (CC BY 4.0), distributed by arXiv under the Creative Commons Attribution 4.0 International licence, http://creativecommons.org/licenses/by/4.0/, as recorded in the arXiv licence metadata for this exact version and shown on the abstract page https://arxiv.org/abs/2604.12340v1. Full licence notice: \texttt{licenses/arxiv-2604.12340-CC-BY-4.0.txt}.\par\smallskip

\noindent\textit{Editorial scope.} Counted unit: the article's proposition prop:phase (source0.tex lines 1302--1343). The article's complete original proof of it is delivered with it (source0.tex lines 1345--1450, one \texttt{proof} environment of 106 lines). No mathematics is written, repaired or re-derived: the statement and the proof are the article's own lines, in the article's own notation, and no retained line is substituted or re-flowed.  Retained uncounted as the source-local context the proof needs: lines 1008--1013; lines 1087--1092; lines 1097--1100; lines 1123--1128; lines 1155--1178.  Nothing was omitted from any retained span, so the package carries no omission record.  No retained line contains a citation, so the wrapper carries no bibliography and no bibliography entry.  No retained line contains a figure environment, an image-inclusion command or drawing code, so no image is delivered and the package declares no asset dependency.  Editorial formatting only, in addition: an AMS \texttt{amsart} preamble that restates the article's own declarations and macros used by the retained lines, namely the environments \texttt{theorem}, \texttt{proposition} on separate counters, together with the standard packages those lines need.  The retained environments print in this document as Theorem 1, Proposition 1 in order of appearance, whereas the article prints the same items with the numbering of its own full document.  The complete native source of the article as distributed in the arXiv e-print for this exact version is bundled unchanged as \texttt{sources/198397/source0.tex}.
\begin{equation}
  2\,\KL\!\big(\mathcal{N}(0,\Sigma_{0}) \,\big\|\, \mathcal{N}(0,\Sigma_{1})\big)
  \;=\;
  \tr(\Sigma_{1}^{-1}\Sigma_{0}) - \NV + \log\frac{|\Sigma_{1}|}{|\Sigma_{0}|}.
  \label{eq:gausKL}
\end{equation}

\begin{equation}
  \frac{2\,\mathrm{Data\ bias}(r)}{\NV}
  \;=\;
  I_{-1}(\lambda_{\mathrm{cut}}) + I_{\log}(\lambda_{\mathrm{cut}}) - r .
  \label{eq:DB-r}
\end{equation}

\begin{equation}
  \mathrm{GE}(r) \;=\; \mathrm{ME}(r) + \mathrm{Data\ bias}(r) ,
  \label{eq:GE-decomp-asymp}
\end{equation}

\begin{equation}
  \frac{d}{dr}\!\left[\frac{2\,\mathrm{GE}}{\NV}\right]
  \;=\;
  f\!\big(\lambda_{\mathrm{cut}}\big) \;-\; f(\epsilon).
  \label{eq:dGEdr}
\end{equation}

\begin{theorem}[Optimal cutoff in $\epsilon$-PCA]
\label{thm:optimal-cutoff}
In the high-dimensional limit, for
$\epsilon \in (\lambda_{-}(\alpha), 1)$, the unique interior
local minimum of the generalization error
$\langle\KL(\PV \,\|\, \QVm)\rangle_{m}$ of
eigen-$\epsilon$-PCA on isotropic Gaussian data with intrinsic
noise floor $\epsilon$ is attained at the cutoff
\begin{equation}
  \boxed{\;\lambda_{\mathrm{cut}}^{*} \;=\; \epsilon\;}
  \label{eq:optcut}
\end{equation}
and the corresponding optimal rank is
\begin{equation}
  \NK^{*} \;=\; \NV \int_{\epsilon}^{\lambda_{+}(\alpha)}
                  p_{\mathrm{MP}}(\lambda;\alpha)\,d\lambda .
  \label{eq:NKopt}
\end{equation}
Equivalently, the optimal model retains exactly those empirical
covariance eigenvalues that exceed the intrinsic noise floor
$\epsilon$. The other root $\widetilde{\lambda} > 1$ of the
first-order condition $f(\lambda) = f(\epsilon)$ is a local
\emph{maximum} of GE on the kept branch.
\end{theorem}

\begin{proposition}[Three-regime phase structure]
\label{prop:phase}
Fix the aspect ratio $\alpha \in (0, 1)$ and consider
$\epsilon$-PCA on isotropic Gaussian data in the
high-dimensional limit $\NV, D \to \infty$ with
$\NV/D = \alpha$. Then the rank ratio
$r^{*} = \NK^{*}/\NV$ that globally minimizes the
generalization error (\ref{eq:GE-decomp-asymp}) satisfies:
\begin{enumerate}
\item[(R1)] (\textbf{retain-all}) \;
  If $\epsilon \leq \lambda_{-}(\alpha) = (1-\sqrt{\alpha})^{2}$,
  then $r^{*} = 1$. Equivalently, the cutoff
  $\lambda_{\mathrm{cut}}^{*}$ lies below the MP support, every
  empirical eigenvalue exceeds $\epsilon$, and the optimal model
  retains all of them.
\item[(R2)] (\textbf{interior}) \;
  If $\lambda_{-}(\alpha) < \epsilon < \epsilon_{*}(\alpha)$,
  then $r^{*}$ is the interior optimum of
  Theorem~\ref{thm:optimal-cutoff},
  $r^{*} = \int_{\epsilon}^{\lambda_{+}(\alpha)}
   p_{\mathrm{MP}}(\lambda;\alpha)\,d\lambda$, and the optimal
  model retains exactly those empirical eigenvalues that exceed
  $\epsilon$.
\item[(R3)] (\textbf{collapse}) \;
  If $\epsilon \geq \epsilon_{*}(\alpha)$, then $r^{*} = 0$. The
  optimal model is the pure noise-floor distribution
  $\mathcal{N}(0,\,\epsilon I_{\NV})$ and uses no information
  from the training data.
\end{enumerate}
The collapse threshold $\epsilon_{*}(\alpha) \in
(\lambda_{-}(\alpha), 1)$ is the unique solution of
\begin{equation}
  \mathrm{GE}_{\mathrm{interior}}(\epsilon)
  \;=\;
  \tfrac{1}{2}\,\NV\!\Big(\tfrac{1}{\epsilon} + \log\epsilon - 1\Big),
  \label{eq:eps-star}
\end{equation}
where the right-hand side is the closed-form value of
$\mathrm{GE}(r{=}0)$ and the left-hand side is the value of
(\ref{eq:GE-decomp-asymp}) at the interior cutoff
$\lambda_{\mathrm{cut}}^{*} = \epsilon$.
\end{proposition}

\begin{proof}
On the half-line $r = 0$ the trained model degenerates to
$\mathcal{N}(0,\, \epsilon I_{\NV})$ and (\ref{eq:gausKL}) gives
$\mathrm{GE}(0) = \tfrac{1}{2}\NV(1/\epsilon + \log\epsilon - 1)$,
the right-hand side of (\ref{eq:eps-star}). On the half-line
$r = 1$ the model is the eigen-MLE and one verifies from
(\ref{eq:dGEdr}) that
$\mathrm{GE}(1) - \mathrm{GE}(\epsilon)$ has the sign of
$\epsilon - \lambda_{-}(\alpha)$ (the MP integrand has support
on $[\lambda_{-},\lambda_{+}]$).

(R1) Suppose $\epsilon \leq \lambda_{-}(\alpha)$. We show that
$f(\lambda) \leq f(\epsilon)$ for every $\lambda \in
[\lambda_{-}, \lambda_{+}]$, so that the right-hand side of
(\ref{eq:dGEdr}) is non-positive on $r \in (0,1)$ and the GE is
non-increasing in $r$, attaining its minimum at $r^{*} = 1$.

Since $f$ is strictly convex on $(0, \infty)$ with global
minimum $f(1) = 1$, the maximum of $f$ on the closed interval
$[\lambda_{-}, \lambda_{+}]$ is attained at one of the
endpoints. For the lower endpoint we have $\epsilon \leq
\lambda_{-} < 1$, and $f$ is strictly decreasing on $(0, 1)$,
so $f(\lambda_{-}) \leq f(\epsilon)$. For the upper endpoint
we use the explicit form $\lambda_{\pm} = (1 \pm \sqrt{\alpha})^{2}$:
writing $x := \sqrt{\alpha} \in (0, 1)$, a direct computation
gives
\begin{equation*}
  f(\lambda_{-}) - f(\lambda_{+})
  \;=\;
  \frac{4x}{(1-x^{2})^{2}}
  \,-\, 2\log\!\frac{1+x}{1-x},
\end{equation*}
which vanishes at $x = 0$ and whose derivative in $x$ is
\begin{equation*}
  \frac{d}{dx}\bigl[f(\lambda_{-}) - f(\lambda_{+})\bigr]
  \;=\;
  \frac{4x^{2}(5 - x^{2})}{(1 - x^{2})^{3}},
\end{equation*}
strictly positive on $(0, 1)$ since the numerator
$4x^{2}(5 - x^{2})$ and the denominator $(1 - x^{2})^{3}$ are
both positive on that interval. Hence
$f(\lambda_{-}) > f(\lambda_{+})$ for every
$\alpha \in (0, 1)$. Combining, $f(\lambda_{+}) <
f(\lambda_{-}) \leq f(\epsilon)$, so $f \leq f(\epsilon)$
throughout the MP support, as claimed.

(R2)--(R3) Suppose $\epsilon > \lambda_{-}(\alpha)$. Then the
interior critical point $\lambda_{\mathrm{cut}}^{*} = \epsilon$
of Theorem~\ref{thm:optimal-cutoff} lies in the MP support and
is a strict local minimum on the open interval. By continuity
of (\ref{eq:GE-decomp-asymp}) on $[0,1]$, the global minimizer
is one of $\{0,\, r_{\epsilon},\, 1\}$, where $r_{\epsilon}$
is the interior critical point. The case $r^{*} = 1$ is
excluded for $\epsilon > \lambda_{-}$ by an analogous endpoint
analysis to that of (R1) above (the trajectory of
$f(\lambda_{\mathrm{cut}}) - f(\epsilon)$ as $r$ approaches $1$
is now strictly positive), so $r^{*} \in \{0, r_{\epsilon}\}$.

We compare $\mathrm{GE}_{\mathrm{int}}(\epsilon) :=
\mathrm{GE}(r_{\epsilon},\epsilon)$ and $\mathrm{GE}(0)$ as
functions of $\epsilon$. Two observations simplify the
computation. First, for fixed $r$, the data-bias term
(\ref{eq:DB-r}) depends on $\epsilon$ only through the
coupling $r = \int_{\lambda_{\mathrm{cut}}}^{\lambda_{+}}
p_{\mathrm{MP}}(\lambda)\,d\lambda$, so at fixed $r$ the cutoff
$\lambda_{\mathrm{cut}}$ and the integrals
$I_{-1}(\lambda_{\mathrm{cut}})$,
$I_{\log}(\lambda_{\mathrm{cut}})$ are all
$\epsilon$-independent, giving
$\partial(\text{Data bias})/\partial\epsilon\big|_{r} = 0$.
Second, by the first-order optimality condition
$\partial\mathrm{GE}/\partial r\big|_{r=r_{\epsilon}} = 0$, the
envelope theorem eliminates the $dr_{\epsilon}/d\epsilon$
contribution. Combining,
\begin{equation*}
  \frac{d\,\mathrm{GE}_{\mathrm{int}}(\epsilon)}{d\epsilon}
  \;=\;
  \frac{\partial \mathrm{ME}}{\partial \epsilon}\bigg|_{r = r_{\epsilon}}
  \;=\;
  \frac{\NV}{2}\,(1 - r_{\epsilon})\,f'(\epsilon),
\end{equation*}
while $\frac{d}{d\epsilon}\mathrm{GE}(0) = \frac{\NV}{2}\,f'(\epsilon)$.
Subtracting,
\begin{equation*}
  \frac{d}{d\epsilon}\big[\mathrm{GE}_{\mathrm{int}}(\epsilon) - \mathrm{GE}(0)\big]
  \;=\;
  -\,\frac{\NV}{2}\,r_{\epsilon}\,f'(\epsilon)
  \;>\; 0
\end{equation*}
on $(\lambda_{-}, 1)$, since $r_{\epsilon} > 0$ and
$f'(\epsilon) = (\epsilon - 1)/\epsilon^{2} < 0$. (Note that
both $\mathrm{GE}_{\mathrm{int}}$ and $\mathrm{GE}(0)$ are
themselves \emph{decreasing} in $\epsilon$; their difference
increases because $\mathrm{GE}(0)$ decreases faster.)
The difference is therefore strictly increasing in $\epsilon$
on $(\lambda_{-}, 1)$, is negative at $\epsilon = \lambda_{-}$
(where $r_{\epsilon} = 1$ and $\mathrm{GE}_{\mathrm{int}}(\lambda_{-}) =
\mathrm{GE}(r{=}1) < \mathrm{GE}(0)|_{\epsilon = \lambda_{-}}$
by the same upper-endpoint calculation as in (R1)), and is
positive as $\epsilon \to 1^{-}$. Hence it changes sign
exactly once at the unique $\epsilon_{*}(\alpha)$ defined by
(\ref{eq:eps-star}). For $\epsilon < \epsilon_{*}$ the interior
wins, giving (R2); for $\epsilon \geq \epsilon_{*}$ the
boundary wins, giving (R3).
\qed
\end{proof}

\end{document}
